\chapter{Groups, cross terms and honest comparisons}\label{ch:groups}
Two actors can share a source without either actor's request being wrong on
its own. The problem is that both requests refer to one source, not two
copies of it. Nonlinear valuation must therefore evaluate their joint change
before asking how any resulting value might be divided.

\section{One group, one endpoint}
For fixed accepted increments $\delta_a$, define
$\delta_G=\sum_{a\in G}\delta_a$. In the cost-free Gaussian domain,
\begin{equation}
 E_G=-\mu(z)^T\delta_G-\tfrac12\delta_G^TH\delta_G.
\end{equation}
Expand the quadratic sum. Symmetry of $H$ gives
\begin{equation}\label{eq:cross}
 E_G=\sum_{a\in G}E_a^{(0)}
       -\sum_{a<b}\delta_a^TH\delta_b,
\end{equation}
where $E_a^{(0)}=V(z)-V(z+\delta_a)$ is the hypothetical value of child
$a$ alone at the same baseline. These singleton values are not sequential
receipts. The cross terms are exactly what their naive sum leaves out.

\section{A shared-source calculation}
At $(18,2)$ with the usual reference and scales, let two children each send
two units along the same edge. Each alone has value seven. Together they
send four and have value twelve, not fourteen. Their increments are both
$(-2,2)$, so $\delta_A^TH\delta_B=2$ and the correction is $-2$.

The same issue occurs when the destinations differ. If two children share
only a source with quantities $a,b>0$, their cross term contains
$ab/\sigma_i^2>0$. The joint field value is smaller than the sum of
their separate source-relief claims by that amount. Disjoint changed
coordinates have zero cross term for the diagonal model.

Opposing increments can instead produce a positive correction. Take
$\delta_A=(-2,2)$ and $\delta_B=(2,-2)$ at the same baseline. Their
standalone values are seven and minus nine. The net increment is zero,
so the group value is zero and the correction is $+2$. Positive
nonadditivity here means cancellation of represented displacement. It is
not proof that real physical circulation is free or socially beneficial.

\section{Same-baseline nonadditivity is not a parallel advantage}
Define $N_G=E_G-\sum_aE_a^{(0)}$. Now choose an admissible sequential
schedule $\pi$, with complete endpoint $x_\pi$. Its field total is
$V(z)-V(x_\pi)$. The comparator-relative difference is
\begin{equation}
 I_{G\mid\pi}=E_G-E_\pi=V(x_\pi)-V(z+\delta_G).
\end{equation}
When fixed additive increments can be serialized feasibly, the final
endpoint is the same and $I_{G\mid\pi}=0$. Yet $N_G$ may be nonzero.
The shared two-unit example has $N_G=-2$ and $I=0$: it contains nonlinear
same-baseline double counting but no endpoint advantage of simultaneity.

A resolver replay can change this comparison. If it recomputes accepted
quantities after the first action, it may produce another endpoint. That
is a different, explicitly named comparator, not evidence that the same
accepted vector has two contradictory values. The bridge source
\cite{bridge} distinguishes quantity-fixed from rule-replayed comparisons.

\section{The conditional Boolean-lattice calculation}
For a set $A$ of fixed children, let
$F(A)=V(z)-V(z+\sum_{a\in A}\delta_a)$ and $F(\varnothing)=0$.
The pair inclusion--exclusion coefficient is
\begin{equation}
 m_{ab}=F(\{a,b\})-F(\{a\})-F(\{b\})
       =-\delta_a^TH\delta_b.
\end{equation}
Because the indicator expansion of $F$ has degree at most two, all
coefficients of order three and higher vanish. This statement needs the
same baseline, potential and fixed increments in every subset.

If a subset reruns a resolver, its increments can differ. The resulting
set function need not be quadratic in membership indicators. Furthermore,
physically inadmissible subsets cannot be turned into valid interventions
by algebraic convenience. Every relevant subset must be admitted under the
declared protocol before a physical Boolean-lattice interpretation is made.
Möbius analysis remains an optional diagnostic, not extra issuance.

\section*{Questions to ask of a group claim}
Were the children valued at one baseline or consecutive states? Were
quantities held fixed or re-resolved? Is the comparator physically admissible?
Does it include the same costs and complete state? Is a correction being
mistaken for a causal contribution? These questions determine the meaning
of an interaction number before its sign is interpreted.

\section*{Chapter recap}
The quadratic cross term explains why independent same-baseline values fail
to close on a shared group. A named sequential comparator answers another
question. Neither calculation supplies an ownership rule.
